\section{Reconstruction at finite-place panels}
\label{sec:finite}

The exact residue link relates the two endpoint observations at each
source label.  When the target panel is nonarchimedean, we show that these
relations force each endpoint family to be a commuting law of homogeneous
analytic boundary maps.  The same maps then describe translated links and
give the upper-speed estimate.  The proof uses compact-open fixed algebras
and Lie algebras over the prime local field; it will also show that the
source panel has the same residue characteristic as the target.

\subsection{The exact link and the frame algebra}

Fix a mixed lexicographic face path missing one type in a non-isolated
component.  The same
construction will subsequently be used with one missing type in each of
several components.  Write $B_S,B_T$ for the source and target panel
boundaries.  Let $S$ be the effective root-generated source panel group,
and let $T$ be the effective target adjoint point group, including
the finite-index overgroup needed for changes of frame in
Convention~\ref{conv:setup-panels}.
In this section $T$ is a group over a finite extension of a prime local
field.  Its prime-field Lie algebra is simple.  The group $S$ is simple
after its ineffective center has been removed, and $T$ embeds faithfully
in the automorphism group of its prime-field Lie algebra.  In particular,
the centralizer in $T$ of an open subgroup of its root-generated group is
trivial.

We first retain the whole ordinary endpoint Hilbert spaces, denoted by
$H'$ at the start and $I$ at the end.  They include the genuine source
amplification of Lemma~\ref{lem:links-genuine-source}, when that amplification is
needed.  The sharp endpoint-face observations and their decompositions
are retained.  Separate compact frames, chosen measurably from these
observations, identify the two target residues with $B_T$.
For $z\in B_S$, the missing-vertex observations in these frames are PVMs
$P_z$ on $H'$ and $Q_z$ on $I$.  We have an isometry
\begin{equation}
 J:\overline I\otimes H'\longrightarrow\mathcal K
 \label{eq:finite-link}
\end{equation}
and an output projectivity PVM $E$ on $T$.  Equation~\eqref{eq:input-6}
says precisely that
\[
 E(dg)J\bigl(\overline{Q_z}(dq)\otimes P_z(du)\bigr)
 \quad\hbox{is supported on }u=gq.
\]
We use here its exact ordered-test meaning established in
Section~\ref{sec:transfer}: clopen graph substitutions are valid on
compact projectivity windows, with bounded output cuts retained.
This is stronger than a relation between the supports of two marginals.

There is an actual source Levi cover $\widetilde S\to S$ acting strongly
continuously on the ordinary endpoint spaces.  It commutes with the
sharp endpoint-face observations, fixes their frame calculi, and moves
the residue labels by $z\mapsto sz$.  Its ineffective finite center
therefore acts trivially on the algebras generated by the residue
observations, even if it acts nontrivially on the Hilbert spaces.
Only this algebra action is asserted to descend to $S$.
All these properties, including the exact input reindexing before a
further limit, are properties of the link already proved in
Sections~\ref{sec:links} and~\ref{sec:transfer}.

Our goal is to recover a fixed homogeneous analytic homeomorphism
$\tau:B_S\to B_T$ such that, on their respective spectral spaces, the
endpoint families are evaluations of maps
\[
 z\longmapsto a\tau(z),\qquad z\longmapsto b\tau(z),
\]
with measurable target projectivity values $a,b$.  The graph equation will
then identify the output projectivity as $ab^{-1}$.  Both commutation of
the endpoint families and this homogeneous form are conclusions to be
proved.

We put all external target frames into one algebra.  For a locally
constant scalar function $f$ on the compact space $B_T$, define the paired
decomposable operator
\begin{equation}
 F_{z,f}(t)=\bigl(f(tP_z),f(tQ_z)\bigr),\qquad t\in T,
 \label{eq:finite-generators}
\end{equation}
on $L^2(T;H')\oplus L^2(T;I)$.  The notation means spectral evaluation:
$f(tP_z)=\int f(tu)P_z(du)$, and likewise at the end.
Let $D_*$ be the von Neumann algebra generated by these operators.
A countable dense set of source labels and countable determining clopen
tests suffice, since the original observations are continuous in the
label on each fixed test.  Thus $D_*$ has separable predual.

In the intended ordinary-map model, the start test in
Equation~\eqref{eq:finite-generators} would be $f(ta\tau(z))$; its frame
and projectivity variables occur only through $x=ta$.  We will prove
that the actual paired algebra has precisely one such intrinsic
$T$-coordinate: $D_*$ is $L^\infty(T)$ with the frame action given by
left translation.  This will establish commutation and reduce the
homogeneous form of the endpoint maps to source equivariance.

Left frame translation defines a continuous action $\alpha$ of $T$:
\[
 (\alpha_a d)(t)=d(a^{-1}t),\qquad
 \alpha_a(F_{z,f})=F_{z,f\circ a^{-1}}.
\]
The actual source transport defines a continuous action $\beta$ of $S$
on $D_*$ with
\[
 \beta_s(F_{z,f})=F_{sz,f}.
\]
These actions commute.  Continuity is in the point-ultraweak topology;
the spatial implementations and boundedness of the generators also
give the corresponding strong continuity on fixed vectors.

The first step is to understand the fixed algebra of a compact open
subgroup of $T$.  The graph equation produces Hilbert--Schmidt
intertwiners between its frame evaluations.  The nonzero eigenspaces
of their compact positive products supply enough finite-dimensional
representations to give the following atomic decomposition.

\begin{proposition}[Compact-open fixed algebras]
\label{prop:finite-frame}
For every compact open subgroup $U<T$, the algebra $D_*^U$ is a
countable product of finite matrix algebras.  Moreover,
\begin{equation}
 D_*^T=\mathbb C1,
 \label{eq:finite-fixed}
\end{equation}
and the source action $\beta$ is nontrivial.
\end{proposition}

\begin{proof}
A decomposable operator fixed by left translation by $U$ is essentially
constant on each left coset.  Thus there are normal evaluations
\[
 \pi_t^+(d)=d_t^+\in B(H'),\qquad
 \pi_r^-(d)=d_r^-\in B(I),\qquad Ut,Ur\in U\backslash T,
 \quad d\in D_*^U.
\]
They are restrictions of the normal coordinate evaluations on
\[
 \ell^\infty\bigl(U\backslash T;
                  B(H')\oplus B(I)\bigr).
\]
They are jointly faithful, and $U\backslash T$ is countable.

Identify $\overline I\otimes H'$ with $\operatorname{HS}(I,H')$ by
$\overline\xi\otimes\eta\mapsto |\eta\rangle\langle\xi|$.
For fixed $Ut,Ur$, put
\[
 \Omega_{t,r}=\{g\in T:Utg=Ur\}.
\]
For every $\zeta\in\mathcal K$, the Hilbert--Schmidt operator
$K=J^*E(\Omega_{t,r})\zeta$ satisfies
\begin{equation}
 d_t^+K=Kd_r^-,\qquad d\in D_*^U.
 \label{eq:finite-hs-intertwining}
\end{equation}

Here is the word-transfer justification of
Equation~\eqref{eq:finite-hs-intertwining}.  First take a word in
generators with real-valued clopen tests and average that word over
$U$.  At frame $u_0t$, the adjoint of the exact graph identity transfers
each start test to the corresponding end test at frame $u_0tg$.
On a compact $g$-window, take a finite partition into relatively clopen
pieces such that each end function $q\mapsto f(u_0tgq)$ in this finite
list is independent of $g$ within a piece, uniformly for $u_0\in U$.
Indeed the action and the tests are
continuous on the compact set of parameters, while a clopen test has
finite image.  Within a piece every factor therefore intertwines
simultaneously.  Transferring a word successively gives the reversed
order on the conjugate Hilbert space, which under the Hilbert--Schmidt
identification is right multiplication by the end word in its ordinary
order.  All the output cuts remain in place.

On $\Omega_{t,r}$, write $tg=u_g r$ with $u_g\in U$.  Haar averaging in
$u_0$ absorbs $u_g$, so the averaged end word has value $d_r^-$ and the
averaged start word has value $d_t^+$.  This proves
Equation~\eqref{eq:finite-hs-intertwining} for averaged words on compact
windows.  Linear combinations and adjoints treat complex tests.
Kaplansky density approximates a bounded element of $D_*$ by bounded
generator polynomials, with arbitrarily small norm slack.  Normal Haar
averaging approximates every element of $D_*^U$ ultraweakly.
The evaluations on this fixed algebra are normal, and multiplication
by a fixed Hilbert--Schmidt operator is weakly continuous on bounded
sets.  Thus the identity extends to all $d\in D_*^U$.
Exhausting $T$ by compact windows removes the temporary localization.
The endpoint frames cause no extra evaluation step: their sharp
decompositions are retained throughout this calculation.

For fixed $t$, the sets $\Omega_{t,r}$ partition the output variable as
$Ur$ varies.  Since $J$ is an isometry, $J^*$ maps onto
$\operatorname{HS}(I,H')$.  The linear span of all the resulting $K$
is therefore dense in Hilbert--Schmidt norm.  The same statement holds
with $r$ fixed and $t$ varying, because right multiplication by a fixed
$g$ permutes the left cosets.
Equation~\eqref{eq:finite-hs-intertwining} and its adjoint give
\[
 KK^*\in\pi_t^+(D_*^U)',\qquad
 K^*K\in\pi_r^-(D_*^U)'.
\]
The nonzero eigenspaces of these compact positive operators are
finite-dimensional reducing spaces for the indicated representations.
Together they span the ordinary evaluation spaces.  For example, a
vector orthogonal to all the ranges of the $K$ would be orthogonal to
the ranges of a Hilbert--Schmidt dense family, and hence would vanish.

Restriction to these reducing spaces supplies a jointly faithful
family of normal finite-dimensional representations of $D_*^U$.
The central support of each such representation is a finite sum of
central atoms, and its algebra on those atoms is a finite sum of finite
matrix algebras.  Joint faithfulness makes the join of these supports
equal to one.  Separability of the predual makes the resulting product
countable.  This proves the first assertion without any restriction
on the endpoint Hilbert-space multiplicities.

If $d\in D_*^T$, its two evaluations are independent of $t,r$.
Equation~\eqref{eq:finite-hs-intertwining} then holds across every
Hilbert--Schmidt operator.  If $AK=KB$ for every rank-one map
$K:I\to H'$ on nonzero Hilbert spaces, testing
$K=|\eta\rangle\langle\xi|$ gives $A=\lambda1$ and $B=\lambda1$
with the same scalar.  This proves Equation~\eqref{eq:finite-fixed}.

Finally suppose that $S$ acted trivially on $D_*$.  Source covariance
and transitivity on $B_S$ would make $F_{z,f}$ independent of $z$.
For countably many labels and tests, evaluate these explicit
decomposable identities at one common generic external frame; the
transformed tests still generate $C(B_T)$.  Continuity in the source
label then makes the original framed $P_z$ independent of $z$ on the
whole start space.  The actual source Levi cover commutes with the
sharp base observation and its frame calculus, so unframing gives the
same statement for the original missing vertex.

Choose a chamber completing the source face, fixing completions in
the other components.  Its full observation is the commuting product
of its canonical vertex observations.  The retained vertices are
fixed by the actual source Levi factor, and the missing vertex has
just been shown fixed.  Hence the full chamber observation is fixed
by this Levi factor as well as by the minimal source parabolic.
These groups generate the larger parabolic allowing the omitted
simple root.  The resulting unital target-equivariant full-boundary
observation contradicts Lemma~\ref{lem:links-obstruction}.
The assertion concerns the original
observations and so descends through the identity amplification.
In a many-panel link, adjoin the other fixed chamber completions
before applying the same argument.  Thus the source action is
nontrivial in every case used here.
\end{proof}

\begin{corollary}[Separation of archimedean types]
\label{cor:finite-arch-separation}
A connected noncompact simple real source cannot be matched to a
finite-place target panel.  Consequently the type permutation preserves
the archimedean and nonarchimedean subsets of the type list.
\end{corollary}

\begin{proof}
A continuous action of a connected group fixes every central atom of
an atomic von Neumann algebra.  To see the relevant discreteness,
a normal state supported on one atom takes values zero and one on
its distinct translates, so that atom has an open stabilizer.
On each finite matrix block the action is a continuous homomorphism
into a projective unitary group.  A noncompact simple real Lie algebra
has no nonzero homomorphism into a compact Lie algebra: simplicity
would make it injective, and a positive definite invariant inner
product would then make the source Lie algebra compact.  A
homomorphism with zero differential from a connected Lie group is
trivial.  Thus the source would fix every $D_*^U$ pointwise.
As compact open $U$ decreases to the identity, the normal averages
$\int_U\alpha_u(d)\,du$ converge ultraweakly to $d$.  Their ranges
therefore generate $D_*$.  This contradicts
Proposition~\ref{prop:finite-frame}.
The archimedean subset of the finite type list is mapped into itself;
finiteness and bijectivity of the type permutation give the assertion
for both subsets.
\end{proof}

\subsection{Faithful corners and finite block actions}

Both panel groups are now nonarchimedean.  The preceding fixed algebras
may have infinitely many matrix blocks.  Passing to a minimal corner
will instead give a sequence of finite-dimensional fixed algebras.  We
will show that their block stabilizers shrink uniformly; this will
identify the corner with the translation algebra of a compact open
subgroup of $T$.  Recovering the full frame algebra will then require
identifying its center and removing its remaining matrix fibres.

Denote the residue characteristics temporarily by $p_S,p_T$.  Choose
small torsion-free logarithmic compact open subgroups $V<S$ and $U<T$,
respectively pro-$p_S$ and pro-$p_T$, in faithful adjoint matrix
coordinates.  Choose a minimal projection $e\in D_*^U$.
We may shrink $V$ to fix $e$.  In fact an atomic central block has an
open stabilizer, by the normal-state test in the preceding proof.
The image of that compact profinite stabilizer in the automorphism
group of its finite matrix block is both profinite and a closed Lie
subgroup, and hence is finite.  Its open kernel fixes the entire
block.  Choose a still smaller logarithmic subgroup inside that kernel.
Put
\begin{equation}
 D_0=eD_*e,\qquad D_0^U=\mathbb Ce.
 \label{eq:finite-corner}
\end{equation}

\begin{lemma}[Automorphisms fixing a covering family of corners]
\label{lem:finite-corner-automorphisms}
Let $A$ be a von Neumann algebra and let $(r_j)$ be projections with
join one.  The group of automorphisms fixing every $r_jAr_j$
pointwise is abelian.
\end{lemma}

\begin{proof}
An automorphism in question fixes $r_i,r_j$ and their central
supports.  Work on the common central support of these two
projections.  The rectangular space $X=r_iAr_j$ is the full selfdual
equivalence bimodule between its endpoint corners: the ultraweak
spans of $XX^*$ and $X^*X$ are those corners on this support.
The restriction of the automorphism to $X$ is right linear, preserves
the right inner product $x^*y$, and has its inverse as adjoint.
It is also left linear.

Adjointable right-module operators on this equivalence bimodule are
left multipliers by the left corner.  One can verify this directly:
finite sums of the rank-one maps $x\mapsto yz^*x$ approximate its
identity strongly; applying an adjointable operator to these sums
produces left multipliers, and their bounded strong limit remains in
the left corner.  Left linearity forces the resulting multiplier to
be central.  Preservation of the inner product makes it unitary.
Thus the automorphisms act on every $r_iAr_j$ by commuting central
unitary multipliers.  These multipliers are themselves in corners
fixed pointwise by every automorphism in the group, so composition
of the actions on each rectangular space is commutative.

It follows that the commutator $\theta$ of any two such automorphisms
fixes every $r_iAr_j$.  For $a\in A$,
$r_i(\theta(a)-a)r_j=0$ for all $i,j$.  Since the ranges of the $r_j$
span the representation space, this gives $\theta(a)=a$.
\end{proof}

\begin{lemma}[Kernel on the corner]
\label{lem:finite-corner-kernel}
The kernel of the $V$-action on $D_0$ has zero Lie algebra.
\end{lemma}

\begin{proof}
This kernel is a closed normal subgroup of $V$, so its Lie algebra
is an ideal in the simple prime-field algebra $\mathfrak s$.
If it were nonzero, it would be all of $\mathfrak s$, and an open
subgroup $V_0<V$ would fix $D_0$ pointwise.
The source and target actions commute, so $V_0$ would also fix
$\alpha_t(e)D_*\alpha_t(e)$ pointwise for every $t\in T$.
The join of the projections $\alpha_t(e)$ is nonzero and $T$-fixed;
it is one by Proposition~\ref{prop:finite-frame}.  A countable
subfamily has the same join.
Lemma~\ref{lem:finite-corner-automorphisms} now shows that every
commutator from $V_0$ acts trivially on $D_*$.  There is a nonidentity
such commutator because $\Lie(V_0)=\mathfrak s$ is nonabelian.
The kernel of the $S$-action is normal in the effective simple group
$S$.  It would therefore be all of $S$, contrary to
Proposition~\ref{prop:finite-frame}.
\end{proof}

Choose a decreasing sequence of compact open normal subgroups
$U_n\triangleleft U$ with intersection $\{1\}$, and set
\[
 D_n=D_0^{U_n}.
\]
The union of these increasing algebras generates $D_0$ ultraweakly,
by normal averaging and continuity of the $U$-action.
The finite group $U/U_n$ acts ergodically on $D_n$.
For a nonzero projection $a\in D_n$, its average is $\lambda_a e$ and
\[
 a\leq [U:U_n]\lambda_a e,\qquad \lambda_a\geq [U:U_n]^{-1}.
\]
An orthogonal family of nonzero projections consequently has at
most $[U:U_n]$ members.  A von Neumann algebra with such a bound is
finite-dimensional: otherwise repeated splitting of nonminimal
projections or its infinite family of atoms gives arbitrarily large
orthogonal families.  Thus every $D_n$ is finite-dimensional.

\begin{lemma}[Block stabilizers]
\label{lem:finite-blocks}
The simple blocks of $D_n$ form a transitive $U$-set.  The stabilizer
$L_n<U$ of a chosen block acts on it by an irreducible projective
unitary representation.  If $G_n<V\times U$ is the stabilizer of the
same block for the commuting source and target actions, then $G_n$
is closed, its projection onto $V$ is surjective, and
\begin{equation}
 [K_n,K_n]\text{ acts trivially on }D_n,\qquad
 K_n=\{v\in V:(v,1)\in G_n\}.
 \label{eq:finite-block-commutators}
\end{equation}
\end{lemma}

\begin{proof}
A union of block orbits would give an invariant central projection,
so ergodicity gives transitivity.  An element of the chosen block
fixed by $L_n$ extends, by $U$-transport, to an invariant element of
$D_n$.  Therefore the fixed algebra of $L_n$ on that block is
scalar.  In a matrix algebra its automorphism action is implemented
projectively, and this scalar-commutant condition is precisely
irreducibility.  The action on the finite set of blocks is
continuous, so $G_n$ is closed.  For each $v\in V$, transitivity
supplies $u\in U$ carrying the block moved by $v$ back to the chosen
one.  Hence $G_n\to V$ is onto.

An element of $K_n$ fixes every block label: its permutation commutes
with the transitive $U$-action and fixes one label.  On the chosen
block its projective implementer centralizes the projective
representation of $L_n$.  Such a projective centralizer is abelian.
Indeed, for an implementer $A$ and projective representation
implementers $W_l$, write
\[
 AW_lA^*=\chi_A(l)W_l.
\]
The representation multipliers cancel, making $\chi_A$ a scalar
character.  The map $[A]\mapsto\chi_A$ is a homomorphism into the
abelian character group.  Its kernel is trivial by Schur's lemma.
Thus commutators of $K_n$ act trivially on this block.  Commutation
with $U$ and transitivity give the assertion on all blocks.
\end{proof}

\begin{corollary}[Separation of residue characteristics]
\label{cor:finite-primes}
The source and target residue characteristics are equal.  The type
permutation preserves each residue-characteristic subset of the type
list.
\end{corollary}

\begin{proof}
If $p_S\ne p_T$, a closed subgroup of the product of a pro-$p_S$
group and a pro-$p_T$ group is the product of its projections.
Pass to finite quotients: a subgroup of a finite $p_S$-group times
a finite $p_T$-group has projected groups with no common nontrivial
quotient, so is their product.  The assertion for closed subgroups
follows by inverse limits.  Consequently $V\times1\subset G_n$ for
every $n$.  Lemma~\ref{lem:finite-blocks} makes $[V,V]$ act
trivially on all $D_n$ and hence on $D_0$.
The closure of this commutator subgroup has full Lie algebra
$[\mathfrak s,\mathfrak s]=\mathfrak s$; for example, take the
Lie algebra of the compact abelian quotient by its closure.
This contradicts Lemma~\ref{lem:finite-corner-kernel}.
Applying the assertion at every type proves the last statement.
\end{proof}

\subsection{A uniform logarithm lemma}

From now on both prime-field Lie algebras $\mathfrak s,\mathfrak t$
are over the same $\mathbb Q_p$.  The finite-level block stabilizers
vary with $n$, so a subgroup-dependent logarithm constant would not
suffice.  We prove the needed uniform statement.

\begin{lemma}[Uniform saturation in a logarithm chart]
\label{lem:finite-uniform-log}
Let $W$ be a sufficiently small logarithmic compact open subgroup
of a linear $p$-adic Lie group.  There is an integer $c\geq0$,
depending only on $W$, such that for every closed subgroup $H<W$,
\begin{equation}
 \exp(p^cM_H)\subset H,\qquad
 M_H=\operatorname{span}_{\mathbb Z_p}\{\log h:h\in H\}.
 \label{eq:finite-uniform-log}
\end{equation}
The assertion includes $p=2$ and subgroups of smaller Lie dimension.
\end{lemma}

\begin{proof}
In dimension zero, a sufficiently small $W$ is trivial and $c=0$
suffices.  Assume henceforth that the ambient Lie algebra has positive
dimension.  Choose a lattice $\Lambda$ in that Lie algebra $\mathfrak w$
on which matrix exponential and logarithm are inverse, with
$W=\exp\Lambda$.  Shrink it so that
\[
 (\Ad(w)-1)\Lambda\subset p\Lambda\qquad(w\in W).
\]
Write $N=\dim_{\mathbb Q_p}\mathfrak w$ and $M=M_H$.
As a submodule of the finite free $\mathbb Z_p$-module $\Lambda$,
$M$ is finitely generated, free, and closed in its rational span.
Choose finitely many $h_1,\ldots,h_a\in H$ whose logarithms generate
$M$.  Such a subfamily exists because $\mathbb Z_p$ is Noetherian.

For $h\in H$, conjugation by $h$ and $h^{-1}$ preserves $M$.
Thus $A=\Ad(h)-1$ is an endomorphism of $M$.  Its characteristic
polynomial on the fixed ambient space has the form
\[
 x^N+a_{N-1}x^{N-1}+\cdots+a_0,\qquad a_j\in p\mathbb Z_p.
\]
Cayley--Hamilton applied on $M$, rather than just on the ambient
lattice, gives
\begin{equation}
 A^NM\subset pM,\qquad
 A^kM\subset p^{\lfloor k/N\rfloor}M.
 \label{eq:finite-cayley-lattice}
\end{equation}
Set
\[
 b=\max\left(0,\sup_{k\geq1}
                    \{v_p(k)-\lfloor k/N\rfloor\}\right)<\infty.
\]
The logarithm series and Equation~\eqref{eq:finite-cayley-lattice}
give
\[
 \log(1+A)M\subset p^{-b}M.
\]
In the fixed matrix logarithm chart,
$\log\Ad(h)=\operatorname{ad}(\log h)$.
Extending over the chosen generators of $M$ therefore proves
\[
 [M,M]\subset p^{-b}M.
\]
Choose $c\geq b+2$ and put $L=p^cM$.  Then
$[L,L]\subset p^2L$.

We check the BCH estimate at every prime.  In Dynkin's formula a
homogeneous term of degree $m$ has a denominator of the form
$km\prod_i r_i!s_i!$, where the $k$ blocks are nonempty and
$\sum_i(r_i+s_i)=m$.  Its $p$-valuation is at most
$m-1+v_p(m)$: the factorial terms contribute at most $m-k$, and
$v_p(k)\leq k-1$.  A bracket of length $m$ in $L$ lies in
$p^{2(m-1)}L$.  For $m\geq3$ the resulting term therefore lies in
$p^{m-1-v_p(m)}L\subset pL$, with its exponent tending to infinity.
The degree-two term is $[X,Y]/2\in pL$, including when $p=2$.
Thus BCH converges, preserves $L$, and satisfies
\begin{equation}
 \operatorname{BCH}(X,Y)\equiv X+Y\pmod{pL}.
 \label{eq:finite-bch}
\end{equation}
Every nonlinear term contains both variables.  The same estimates
therefore give, for $r,s\geq0$,
\[
 \operatorname{BCH}(X,Y)-X-Y\in p^{r+s+1}L
 \quad(X\in p^rL, Y\in p^sL).
\]
In particular $\exp L$ and $\exp(p^jL)$ are groups.

At level $j$, the elements $h_i^{p^{c+j}}\in H$ have logarithms
$p^{c+j}\log h_i$ spanning $p^jL/p^{j+1}L$.
Given $x\in L$, choose a product of these powers at level zero
whose logarithm agrees with $x$ modulo $pL$.  Its inverse times
$\exp x$ has logarithm in $pL$.
Correct this error with a product at level one, and continue.
Equation~\eqref{eq:finite-bch} and its refined estimate improve the
error by one level at every step.  The correcting factors tend to
one, and their successive products in $H$ converge to $\exp x$.
Closedness of $H$ proves Equation~\eqref{eq:finite-uniform-log}.
Every constant was chosen before $H$, as required.
\end{proof}

We use the usual analytic structures on closed subgroups and
quotients of $p$-adic Lie groups, and the fact that continuous
homomorphisms between them are analytic; see
\cite[Propositions~2.2--2.3]{Gloeckner2008}.
In particular a surjective homomorphism between compact $p$-adic
Lie groups has surjective differential.  Indeed, the image of a
sufficiently small exponential open subgroup is a compact
finite-index subgroup of the target and is therefore open.
Naturality of exponential puts its logarithm in the differential
image, which consequently spans the target Lie algebra.

\subsection{Hausdorff limits and equality of panel dimensions}

We now use the uniform logarithm estimate to control the block
stabilizers.  The compact metrizable group $V\times U$ has a compact
space of closed subgroups in the Hausdorff topology.  Every limit
$G_\infty$ of the $G_n$ still projects onto $V$: for fixed $v\in V$,
choose $(v,u_n)\in G_n$ and pass to a convergent subsequence of $(u_n)$.
The key point is that a nonzero source-only Lie kernel in such a limit
would force one fixed open source subgroup into all sufficiently late
stabilizers.  Its commutators would then fix the whole corner, contrary
to Lemma~\ref{lem:finite-corner-kernel}.  The next argument applies to
every choice of a block at every level.

\begin{lemma}[The limiting source-only kernel]
\label{lem:finite-limit-kernel}
For any sequence of levels tending to infinity, any choices of
blocks at those levels, and any Hausdorff limit $G_\infty$ of their
stabilizers, the subgroup $G_\infty\cap(V\times1)$ has zero Lie
algebra.
\end{lemma}

\begin{proof}
Its Lie algebra is invariant under $\Ad(V)$, because
$G_\infty\to V$ is onto.  It is thus an ideal of the simple algebra
$\mathfrak s$.  If nonzero it equals $\mathfrak s$, and
$G_\infty$ contains $V_0\times1$ for an open subgroup $V_0<V$.
Set $d=\dim_{\mathbb Q_p}\mathfrak t$.
Choose finitely many $v_1,\ldots,v_a\in V_0$ such that
\begin{equation}
 \sum_{i=1}^a
   \ran\bigl((\Ad(v_i)-1)^d\bigr)=\mathfrak s.
 \label{eq:finite-spanning}
\end{equation}
For this choice, start with a noncentral sufficiently small element
of a split torus inside $V_0$.  Its adjoint action is semisimple
over a splitting field, so the nonzero image of $\Ad(v)-1$ is
unchanged upon taking a positive power.  The span of its conjugates
by $V_0$ is invariant under this open subgroup, hence under
$\mathfrak s$.  It is a nonzero ideal and equals $\mathfrak s$.
Finitely many conjugates suffice and give
Equation~\eqref{eq:finite-spanning}.

Pass to the chosen convergent subsequence, denoting its levels again
by $n$.  Hausdorff convergence gives
\[
 (v_{i,n},w_{i,n})\in G_n,\qquad
 (v_{i,n},w_{i,n})\longrightarrow(v_i,1).
\]
Let
\[
 M_n=\operatorname{span}_{\mathbb Z_p}\log G_n,\qquad
 M_V=\operatorname{span}_{\mathbb Z_p}\log V.
\]
Surjectivity of $G_n\to V$ gives the exact equality
$\operatorname{pr}_{\mathfrak s}M_n=M_V$.
All these modules lie in fixed logarithm lattices.
Put
\[
 \chi_{i,n}(X)=\det\bigl(X1-(\Ad(w_{i,n})-1)|_{\mathfrak t}\bigr),
 \qquad
 B_{i,n}=\chi_{i,n}(\Ad(v_{i,n})-1).
\]
The polynomials $\chi_{i,n}$ have integral coefficients and converge
coefficientwise to $X^d$.  Conjugation by $(v_{i,n},w_{i,n})$
preserves $M_n$.  Therefore
$\chi_{i,n}(\Ad(v_{i,n},w_{i,n})-1)$ preserves $M_n$, and
Cayley--Hamilton annihilates its target component exactly.
For every $x\in M_V$ choose a lift $(x,y)\in M_n$ and apply this
polynomial.  It follows that
\begin{equation}
 \left(\sum_i B_{i,n}M_V\right)\times\{0\}\subset M_n,\qquad
 B_{i,n}\longrightarrow(\Ad(v_i)-1)^d.
 \label{eq:finite-polynomial-lattice}
\end{equation}

Choose a basis of the fixed lattice $M_V$ and concatenate the
matrices $B_{i,n}$.  Equation~\eqref{eq:finite-spanning} gives a
nonzero full-rank square minor in the limiting concatenation.
The corresponding integral determinants at all sufficiently late
levels have one fixed finite valuation, say $a$.
The adjugate identity then yields
\[
 p^aM_V\times\{0\}
 \subset\left(\sum_iB_{i,n}M_V\right)\times\{0\}
 \subset M_n
\]
uniformly at those levels.
Apply Lemma~\ref{lem:finite-uniform-log} in the fixed product
logarithm group $V\times U$.  With a constant $c$ independent of
$n$, it follows that
\[
 \exp(p^{c+a}M_V)\times1\subset G_n.
\]
After increasing $a$ if needed, the displayed source exponential
is a fixed open subgroup $V_1<V$.
Equation~\eqref{eq:finite-block-commutators} says that $[V_1,V_1]$
fixes $D_n$ at every sufficiently late level of this subsequence.
These levels tend to infinity and their algebras exhaust $D_0$.
Thus $[V_1,V_1]$ fixes $D_0$.
Its closure has full Lie algebra
$[\mathfrak s,\mathfrak s]=\mathfrak s$, contradicting
Lemma~\ref{lem:finite-corner-kernel}.
\end{proof}

\begin{proposition}[Dimensions and shrinking stabilizers]
\label{prop:finite-dimensions}
Matched finite-place panel groups have equal Lie dimension over
their common prime field.  The target stabilizers of the blocks of
$D_n$ shrink to the identity uniformly over all choices of blocks.
The corner $D_0$ is abelian and, as a $U$-algebra, is
$L^\infty(U)$ with translation and Haar probability.
\end{proposition}

\begin{proof}
Let $G_\infty$ be any limit.  Its Lie projection onto
$\mathfrak s$ is surjective.  By
Lemma~\ref{lem:finite-limit-kernel}, its Lie projection into
$\mathfrak t$ is injective.  Consequently
\begin{equation}
 \dim_{\mathbb Q_p}\mathfrak s
 \leq\dim_{\mathbb Q_p}\mathfrak t.
 \label{eq:finite-dimension-inequality}
\end{equation}
Use the same missing-root construction at every type.
Source and target panel dimensions come from the same finite list,
and Corollaries~\ref{cor:finite-arch-separation}
and~\ref{cor:finite-primes} keep each type cycle over a single
prime field.  Iterating
Equation~\eqref{eq:finite-dimension-inequality} around that cycle
forces equality at every step.  Only effective prime-field Lie
dimensions enter this argument; no ineffective compact factor or
complex Hilbert-space dimension is counted.

Both projections of $\Lie(G_\infty)$ are now isomorphisms.
The target-only subgroup $G_\infty\cap(1\times U)$ has dimension
zero.  A compact zero-dimensional analytic subgroup is finite,
and $U$ is torsion-free, so
\begin{equation}
 G_\infty\cap(1\times U)=\{1\}.
 \label{eq:finite-target-kernel}
\end{equation}
Suppose a neighborhood of one failed to contain every sufficiently
late target block stabilizer.  Choose a smaller clopen neighborhood,
levels tending to infinity, offending blocks, and stabilizer
elements $u_n$ outside it.  Along a subsequence their $G_n$ converge
to $G_\infty$ and $u_n\to u\ne1$.
Then $(1,u)\in G_\infty$, contradicting
Equation~\eqref{eq:finite-target-kernel}.
This proves uniform shrinking for arbitrary block choices.

Fix $m$.  For all sufficiently large $n$, the stabilizer of every
block of $D_n$ lies in $U_m$.  It therefore fixes the image of
$D_m$ in that block pointwise.  Its projective representation is
irreducible by Lemma~\ref{lem:finite-blocks}, so that image is
scalar.  The block representations are jointly faithful, and
therefore $D_m$ is abelian.  Exhaustion makes $D_0$ abelian.

Write the spectrum of $D_n$ as the finite transitive $U$-space
$X_n$.  The inclusions give surjective equivariant maps
$X_{n+1}\to X_n$.  Averaging a faithful normal state of $D_0$
over $U$ gives a faithful invariant state whose restriction to
each $X_n$ is the uniform probability.  Thus
\[
 D_0=L^\infty(X,\nu),\qquad X=\varprojlim X_n,
\]
with these uniform marginals.  The compact group $U$ acts
transitively on $X$: the sets of elements carrying the first $n$
coordinates of one point to those of another are nonempty nested
closed subsets of $U$.  Their intersection is nonempty.
The stabilizer of a point belongs to all its finite-level
stabilizers, so uniform shrinking makes it trivial.  Hence $X$
is the translation $U$-torsor, and $\nu$ is Haar probability.
\end{proof}

\subsection{The homogeneous center and elimination of matrix fibres}

\begin{proposition}[The full frame algebra]
\label{prop:finite-torsor}
There is a $T$-equivariant normal isomorphism
\[
 D_*\cong L^\infty(T)
\]
with the left translation action.  In these coordinates the
commuting source action has the form
\begin{equation}
 (\beta_s f)(x)=f(x\eta(s)),\qquad
 \eta:S\longrightarrow T,
 \label{eq:finite-right-action}
\end{equation}
where $\eta$ is a continuous homomorphism with isomorphic
$\mathbb Q_p$-differential.
\end{proposition}

\begin{proof}
By Proposition~\ref{prop:finite-dimensions}, $e$ is an abelian
projection.  On its central support $z_e$ the algebra is type I,
and the map
\begin{equation}
 Z(D_*)z_e\longrightarrow eD_*e,\qquad a\longmapsto ae
 \label{eq:finite-center-corner}
\end{equation}
is a normal isomorphism.  In the type I decomposition of this
central summand, $e$ has rank one in every factor, and full central
support makes the scalar field faithful, which proves the assertion.
The map is $U$-equivariant because $U$ fixes $e$.
Countably many $T$-translates of $e$ have join one, so their central
supports cover one and $D_*$ is type I.
Its center has a positive-measure $U$-invariant chart isomorphic
to the translation space $U$.  Its $T$-action is ergodic by
Equation~\eqref{eq:finite-fixed}.

Use the standard Borel point realization for a continuous action
of a locally compact second countable group on a separable measure
algebra~\cite{Mackey1962}.  Let $X$ realize the center, and let
$\varphi:U\to X$ be the measurable isomorphism onto the chart.
Equivariance holds almost everywhere for each group element.
Fubini in the two Haar variables gives a point $x_0\in X$ with
\[
 \varphi(u)=ux_0\quad\hbox{for Haar-almost every }u\in U.
\]
Explicitly, choose a generic $u_0$ such that
$\varphi(vu_0)=v\varphi(u_0)$ for almost every $v$, and put
$x_0=u_0^{-1}\varphi(u_0)$.
Let $L=\operatorname{Stab}_T(x_0)$.
If $h\in L\cap U$, compare $u$ and $uh$ in a common conull
injectivity set of $\varphi$; their images agree, so $h=1$.
A subgroup meeting a neighborhood of one only at one is discrete
and closed.  The orbit map $T/L\to X$ is a Borel injection and
has a Borel inverse on its image.  Countably many translates of
the chart cover the center modulo null sets, so this orbit is
conull.  The measure class is quotient Haar class, as follows on
the chart and its translates.  We have obtained
\[
 Z(D_*)\cong L^\infty(T/L).
\]

The measurable automorphisms of $T/L$ commuting with left
translation are exactly right translations from $N_T(L)/L$.
For such an automorphism $F$, the function
$g\mapsto g^{-1}F(gL)$ is essentially left invariant and is
therefore essentially constant, say $nL$.
Thus $F(gL)=gnL$ almost everywhere.  Right-$L$ consistency gives
$n^{-1}Ln\subset L$, and the same argument applied to $F^{-1}$
gives equality.  Conversely a normalizer coset defines such an
automorphism.  The injection from $N_T(L)/L$ into the
measure-class automorphism group is Borel, with Borel inverse
onto its image.  Indeed its action is measurable on a countable
determining collection of sets, and the Borel inverse follows
from injectivity between standard Borel spaces.
The commuting source action consequently gives a measurable
homomorphism
\[
 \eta:S\longrightarrow N_T(L)/L,
\]
where the point action uses inverse right translation so that
the algebra convention is Equation~\eqref{eq:finite-right-action}.
The normalizer is closed.  Measurable homomorphisms between
locally compact second countable groups are continuous, so
$\eta$ is continuous and analytic over $\mathbb Q_p$.

Its differential is injective.  Otherwise its kernel has
positive Lie dimension; a corresponding small subgroup of $V$
acts trivially on the center.  Since $V$ fixes $e$,
Equation~\eqref{eq:finite-center-corner} would make that subgroup
act trivially on $eD_*e$, contradicting
Lemma~\ref{lem:finite-corner-kernel}.
Because $L$ is discrete and
$\dim\mathfrak s=\dim\mathfrak t$, the closed Lie subgroup
$N_T(L)$ has the full dimension of $T$ and is open.
For each $l\in L$, conjugation from $N_T(L)$ into the discrete
group $L$ is continuous.  The element $l$ therefore centralizes
an open subgroup.  The effective open-centralizer assertion from
Section~\ref{sec:setup} forces $l=1$.
Thus $L=\{1\}$, the center is the $T$-torsor, and $d\eta$ is an
isomorphism.

It remains to remove matrix fibres.  The type I decomposition
and ergodicity of the center give
\[
 D_*\cong L^\infty(T;B(\mathcal V))
\]
with one separable Hilbert space $\mathcal V$ of constant
dimension.  In a measurable trivialization the action is
\[
 (\alpha_a d)(x)=c(a,x)\bigl(d(a^{-1}x)\bigr),\qquad
 c(ab,x)=c(a,x)c(b,a^{-1}x),
\]
where $c$ takes values in $\operatorname{Aut}(B(\mathcal V))$.
No unitary lift of this automorphism cocycle is needed.
Changing variables in its cocycle identity and applying Fubini
allows one to freeze a generic starting point $x_0$ and define
\[
 b(x)=c(xx_0^{-1},x).
\]
The identity becomes
\[
 c(a,x)=b(x)b(a^{-1}x)^{-1}
 \quad\hbox{for almost every }(a,x).
\]
Changing the fibre trivialization by $b$ removes the cocycle.
Continuity of the action extends the equality from almost every
$a$ to every $a$.  Pure translation has fixed algebra the
constant fields $B(\mathcal V)$.  Equation~\eqref{eq:finite-fixed}
therefore forces $\dim\mathcal V=1$.
\end{proof}

\subsection{Ordinary homogeneous panel maps}

\begin{proposition}[Recovery of the original panel observations]
\label{prop:finite-maps}
The family $(P_z)_{z\in B_S}$ commutes on the original start
space, and the family $(Q_z)_{z\in B_S}$ commutes on the original
end space.  There is a homogeneous analytic homeomorphism
$\tau:B_S\to B_T$ and measurable target projectivity observations
$a,b$ on their respective spectra such that these families are
evaluations of
\begin{equation}
 z\longmapsto a\tau(z),\qquad
 z\longmapsto b\tau(z).
 \label{eq:finite-panel-laws}
\end{equation}
In the independent input tensor product, the output projectivity
is sharp and equals $ab^{-1}$:
\begin{equation}
 E(W)J=J\,1_W(ab^{-1})
 \quad\hbox{for every Borel }W\subset T.
 \label{eq:finite-projectivity-sharp}
\end{equation}
The assertions hold coordinatewise on many-panel links.
\end{proposition}

\begin{proof}
The commutativity of $D_*$ makes the commutators of its paired
generators vanish at almost every external frame.  Use a common
conull frame set for countably many labels and clopen tests.
At one such frame the transformed tests still generate
$C(B_T)$, so the original $P_z$ commute, as do the $Q_z$, for
the determining labels.  Continuity in the source label extends
commutation to every fixed pair of labels.

In the intrinsic coordinates $D_*=L^\infty(T)$ of
Proposition~\ref{prop:finite-torsor}, the spectral values of the
generators give a boundary-valued measurable function $F(x,z)$.
Countable determining observations and continuity in measure
give a jointly measurable version.  More explicitly, use a
faithful probability in the Haar class and approximate the compact
label space by finite simple label choices whose errors in measure
are summable.  The almost-everywhere limit of the resulting
jointly measurable approximations represents each fixed label;
the countable boundary tests determine the boundary value.
Frame covariance yields
\[
 F(ax,z)=aF(x,z)
\]
almost everywhere, and Fubini gives $F(x,z)=xF_0(z)$.
The commuting source action then gives
\begin{equation}
 F_0(sz)=\eta(s)F_0(z)
 \label{eq:finite-equivariant-map}
\end{equation}
in the ordinary measure class on the source boundary.

We explain why $d\eta$ identifies the original boundaries.
Regard each effective adjoint group over its coefficient field
as its restriction of scalars to $\mathbb Q_p$.
The resulting Lie algebra is the prime-field Lie algebra used
above.  In characteristic zero, all derivations of a semisimple
Lie algebra are inner, and the identity component of its
algebraic automorphism group is the corresponding adjoint
algebraic group.  Conjugation by $A=d\eta$ therefore gives an
isomorphism of these adjoint algebraic groups over $\mathbb Q_p$.
Differentiating the homomorphism identity gives
\begin{equation}
 A\Ad(s)=\Ad(\eta(s))A,\qquad s\in S.
 \label{eq:finite-adjoint-intertwining}
\end{equation}
In faithful adjoint coordinates the algebraic isomorphism on
the root-generated point groups is consequently exactly $\eta$.

An algebraic group isomorphism carries rational parabolics and
their unipotent radicals to the corresponding groups.
The rational parabolics of a restriction of scalars are precisely
the restrictions of the original rational parabolics.
For completeness, after a finite separable splitting extension
the restriction of scalars is a product indexed by the
embeddings of the coefficient field.  A parabolic of a product
is a product of parabolics; descent permutes the conjugate
factors.  Choosing one factor identifies the descent condition
with a parabolic over the original field.  The same argument
applies if the effective almost-simple group itself has several
absolutely simple factors in one Galois orbit.  Minimal rational
parabolics thus recover the original rank-one boundary, with
no splitness assumption and no prior equality of coefficient
fields.  This produces a homogeneous analytic homeomorphism
$\tau:B_S\to B_T$ intertwining $S$ with $\eta(S)$.

Write $B_S=S/P$.  By
Equation~\eqref{eq:finite-equivariant-map}, the measurable lift
\[
 s\longmapsto\eta(s)^{-1}F_0(sP)
\]
is essentially left invariant, hence essentially constant.
Right-$P$ consistency makes this constant fixed by $\eta(P)$.
A minimal parabolic fixes exactly its own point on a rank-one
boundary: its unipotent radical fixes that point and acts
transitively on the complementary open Bruhat cell, which has
more than one point.  Therefore this constant is the parabolic
point defining $\tau$, and $F_0=\tau$ almost everywhere.
Both the original generators and the functions
$x\mapsto f(x\tau(z))$ are continuous in measure in $z$.
The generator identity consequently holds in $D_*$ for every
fixed label, not only for almost every label.

We now remove the external frame without using point evaluation
on arbitrary elements of a diffuse $L^\infty$ algebra.
Let $\Omega^+$ be a standard spectrum of the commutative algebra
of the original start observations, with a faithful normal
probability class.  The normal representation of
$D_*=L^\infty(T)$ by its start decomposable fields takes values
in $L^\infty(T_{\mathrm{ext}}\times\Omega^+)$.
It is pullback by a measurable map $h(t,\omega)$ into the
intrinsic $T$ variable.  The explicit generator identities read
\[
 t\,p_z(\omega)=h(t,\omega)\tau(z)
\]
for almost every $(t,\omega)$, simultaneously for countably
many labels and tests.  Choose one generic external frame $t_0$
and put $a(\omega)=t_0^{-1}h(t_0,\omega)$.
Then $p_z(\omega)=a(\omega)\tau(z)$ for the determining labels.
Each map $a(\omega)\tau$ is continuous on the compact source
boundary.  Bounded convergence of its spectral evaluations and
the original label continuity extend the equality to every
fixed label.  The end representation gives $b$ in the same way.
Only these countably many measurable generator identities were
evaluated at $t_0$; no normal evaluation on all of
$L^\infty(T)$ is asserted.

On the independent input tensor product, $a$ and $b$ commute.
Substitution of Equation~\eqref{eq:finite-panel-laws} in the
exact graph link gives
\[
 a\tau(z)=g\,b\tau(z)
\]
for countably densely many labels simultaneously.  Products of
the input spectral restrictions for these labels are now
legitimate, since all these input observations commute.
The map $\tau$ is onto and the target action is faithful, so
the common zero set of these equations is exactly $g=ab^{-1}$.
On compact sets in $T^3$, finite clopen partitions and successively
more determining labels give the corresponding graph
substitution.  Exhausting the compact sets and then using normal
spectral calculus gives
Equation~\eqref{eq:finite-projectivity-sharp}.
Equivalently, for a clopen output cut its range annihilates the
complementary input graph cut and conversely; the two identities
give the displayed intertwining, and monotone approximation
extends it to Borel cuts.

The map from target projectivities to continuous homogeneous
boundary maps is continuous and injective, and its inverse onto
its image is Borel.  Thus the recovered homeomorphism
observations and their inverses are measurable in the usual
map spaces.  The argument uses just the selected coordinate of
the link.  Other measured coordinates remain in the coefficient
spaces, so the same conclusion holds separately on each
coordinate of a many-panel link.  Commutation between different
coordinates or different measured copies will be established in
Section~\ref{sec:gluing}; it has not been used here.
\end{proof}

\subsection{Exact translated laws and uniform upper speed}

For each finite-place effective panel group $R$ in the fixed type
list, choose a norm on its prime-field adjoint Lie algebra, with
$|p|=p^{-1}$ on $\mathbb Q_p$.  Let $\ell_R$ be its nonnegative
rank-one Cartan gap in the module-absolute-value convention of
Section~\ref{sec:setup}.  There is a constant $d_R>0$, depending
only on this rooted panel type, such that
\begin{equation}
 \log\|\Ad(r)\|=d_R\ell_R(r)+O(1),\qquad r\in R,
 \label{eq:finite-adjoint-growth}
\end{equation}
with uniform error.  On a dominant split translation the largest
adjoint eigenvalue has logarithm a fixed positive multiple of the
single simple-root coordinate.  The conversion from the module
absolute value over the coefficient field to the prime-field
norm changes this fixed multiple. More explicitly, if $k$ is the
coefficient field and $r_{\max}\in\{1,2\}$ is the largest relative
root multiple in this rank-one group, then
\[ d_R=\frac{r_{\max}}{[k:\mathbb Q_p]}. \]
Indeed $|a|_k=|a|_p^{[k:\mathbb Q_p]}$. This prime-field
conversion is distinct from the $k$-root multiplicities in the Haar
character $2\rho$. Bounded centralizer parts and
the compact Cartan factors change operator norms by bounded
factors, proving Equation~\eqref{eq:finite-adjoint-growth} for
all representatives.  The symmetry of the adjoint roots makes
$d_R$ unchanged by inversion or opposition.  These constants
belong to one list on source and target types.

\begin{proposition}[Uniform finite-place upper speed]
\label{prop:finite-speed}
After the inner limit of a translated panel link, let $t\geq0$
be the total missing-root coordinate of the source translation,
including all slower offsets in that coordinate.  The translation
is taken on the permissible cocharacter mesh of
Section~\ref{sec:setup}.  If $g$ is the target residual
projectivity, then
\begin{equation}
 d_T\ell_T(g)\leq d_St+O_{\mathrm{tight}}(1).
 \label{eq:finite-speed}
\end{equation}
Here tightness is on each fixed ordinary input vector, uniformly
in the slower translations with this total coordinate as center.
The assertion holds coordinatewise, and simultaneously as an
upper-tail estimate, on many-panel links.
\end{proposition}

\begin{proof}
Use the convention in which the translated start observation is
$P_z^{(t)}=P_{s_tz}$ on its original input space and the end observation
is $Q_z$.  Thus $s_t$ is the effective start-label transport; it may be
the inverse of the actual completing displacement, and
$\ell_S(s_t)=t$ in either orientation.
The permissible mesh ensures that $s_t$ belongs to the
root-generated group $S$, so $\eta(s_t)$ is defined.
The exact translated-input identity of
Proposition~\ref{prop:transfer-residue-graph} and
Proposition~\ref{prop:finite-maps} now give
\[
 a\tau(s_tz)=g\,b\tau(z).
\]
Since $\tau(s_tz)=\eta(s_t)\tau(z)$, this identifies
\begin{equation}
 g=a\eta(s_t)b^{-1},
 \label{eq:finite-translated-law}
\end{equation}
where $s_t$ is deterministic with total missing gap $t$.
The Hilbert-space transport uses the actual source element or
its actual cover; only its effective label action occurs in
Equation~\eqref{eq:finite-translated-law}.

Crucially, $a$ and $b$ in this equation are the observations on
the original independent input spaces.  Source covariance moves
the entire translation dependence into the middle factor.
All slower translations with nonzero missing coordinate have
already been combined in $s_t$; replacing a completing time
$t_0$ by $t_0+u$ replaces this factor by $\eta(s_{t_0+u})$.
Other split translations have zero coordinate in this residue
and act trivially on its observations.  Any remaining bounded
effective actions range in a fixed compact set and can be
included in the compact multipliers.  The sharp endpoint bases
and their compact frame functions retain their fixed input laws
throughout these translations.  These statements hold after
the inner limit for each fixed choice of the slower parameters.

With $A=d\eta$, Equation~\eqref{eq:finite-adjoint-intertwining}
gives
\[
 \Ad(\eta(s_t))=A\Ad(s_t)A^{-1},\qquad
 \log\|\Ad(\eta(s_t))\|\leq d_St+C.
\]
The same estimate holds for $s_t^{-1}$.
Fix an input vector $\xi\in\overline I\otimes H'$ and
$\varepsilon>0$.  The commuting map observations $a,b$ have
a joint scalar spectral measure on this vector.
Choose compact sets containing both observations except for
squared input norm at most $\varepsilon\|\xi\|^2$.
On that joint cut, $\|\Ad(a)\|$ and $\|\Ad(b^{-1})\|$ have
fixed bounds independent of $t$ and all slower parameters.
Submultiplicativity, Equation~\eqref{eq:finite-translated-law},
and Equation~\eqref{eq:finite-adjoint-growth} yield
\[
 d_T\ell_T(g)\leq d_St+C_{\xi,\varepsilon}
\]
on this cut, with a constant independent of those parameters.
The exact graph intertwining and the isometry of the link show
that the excluded output squared norm is at most the excluded
input squared norm.  This proves
Equation~\eqref{eq:finite-speed}, including for input vectors
that are not simple tensors.

For finitely many panel coordinates, apply the same argument
to each coordinate's ordinary input laws and sum their
exceptional output masses; no commutation of the different
input map systems is needed for this upper-tail union bound.
All constants are independent of the slower parameters.
Consequently the estimate persists through further prescribed
iterated limits with nonnegative total coordinate centers.
\end{proof}

At a finite-place non-isolated type $i$, define the comparison
constant $s_i=d_S/d_T$.  Opposition leaves it unchanged, and
its product around every type cycle is one, because the $d_R$
are drawn from a single type list.  Proposition~\ref{prop:finite-speed}
is the finite-place upper-speed input for the simultaneous height
calculation.  The isolated types left by
Corollary~\ref{cor:finite-arch-separation} are archimedean and
will be treated by the separate estimates below.
